% === VII. INTEGRATION WITH THE CONTROLLER ====================================
\section{Integration with the Controller}
\label{sec:integration}

The identified model leaves the pipeline through one interface. A nonlinear MPC
path-tracking controller (acados~\cite{acados2022}, partial-condensing
HPIPM~\cite{hpipm2020}, horizon $N=20$ at a \SI{20}{\hertz} control period)
carries the same single-track model as the identifier and consumes $\phip$
directly through \eqref{eq:mf}, pushed over a ROS~2 service at the end of each
identification cycle. We claim no contribution to the controller design; it
appears here as the consumer that defines what an admissible identification is.

\subsection{Admissibility gates}
\label{subsec:gates}

A coefficient set can lie strictly inside every box constraint and still be
unusable. We observed two ways, each requiring a gate on a \emph{derived}
quantity rather than a coefficient:

\begin{enumerate}
\item \textbf{Grip ceiling versus trajectory demand.} The controller computes
$\max_{k}(v_{x,k}^{2}|\kappa_{k}|)$ over the speed-clamped reference and
rejects any identification whose $a_{y}^{\max}$ from \eqref{eq:gripceiling}
falls below it. On the earlier trajectory the degenerate fit reported a
\SI{0.40}{\g} ceiling against a \SI{1.22}{\g} demand and produced
\SI{41.89}{\meter} of cross-track error at \SI{27}{\meter\per\second},
against \SI{0.03}{\meter} for a physically plausible set ($D=\num{1.5}$) on the
same trajectory, at the same speed and through the same controller; the same
identified set gives \SI{0.16}{\meter} at \SI{11.1}{\meter\per\second} and
\SI{0.43}{\meter} at \SI{15}{\meter\per\second}, the $v^{2}$ scaling the
Introduction describes. These figures come from an offline closed-loop harness rather than a controlled
pair of simulator runs, and are quoted as the observation that motivated the
gate, not as a result. The gate is one-sided by construction and catches only a
ceiling that is too low, so \basearm{}'s \SI{2.00}{\g}
(Section~\ref{subsec:resultstire}) passes it.
\item \textbf{Understeer sign and critical speed.} A set with
$l_{f}C_{f}>l_{r}C_{r}$ whose $v_{\mathrm{crit}}$ falls inside the controller's
speed range makes the prediction model open-loop unstable, and the resulting
$\rho(A_{d})>1$ is raised to $\rho^{N}$ in the condensed QP. An observed set
gave $C_{f}=\SI{37031}{\newton\per\radian}$ against
$C_{r}=\SI{11691}{\newton\per\radian}$, i.e.\
$v_{\mathrm{crit}}=\SI{14.5}{\meter\per\second}$ at the measured geometry and
mass of Table~\ref{tab:vehicle}, inside the controller's speed range on every
trajectory in this paper.
\end{enumerate}

A third observation is a controller-implementation detail rather than a gate on
the identification: an admissible set can still be stiff relative to the control
period, so the prediction model sizes its RK4 sub-step from its own fastest
lateral mode, restoring $\rho\approx1.000$ across the speed range, at a median
QP setup-and-solve cost of \SI{0.36}{\milli\second} against the
\SI{50}{\milli\second} control period.


An identification that passes the gates is a model the controller can be held
responsible for; one that fails is reported as unidentified rather than
deployed. Two further controller-side changes were required before any
identified model could be evaluated at full-scale speeds: reference speeds are
clamped at ingest, and the solved state is rolled forward by the measured
transport delay before each solve. Both are recorded for completeness and are
not identification results.
